\begin{figure}[tb]
\centering
\begin{tikzpicture}[x=1.4cm,y=1.05cm,>=Stealth,font=\small]
  \draw[->] (-2.4,0)--(3.05,0) node[right] {$y$};
  \draw[->] (-2.25,-.1)--(-2.25,3.55) node[above] {$u$};
  \draw[densely dashed,gray] (-2.25,.85)--(2.65,.85);
  \node[left] at (-2.25,.85) {$h$};
  \draw[thick] (-1.65,.85) .. controls (-1.15,2.50) and (-.45,3.0) .. (.30,3.0)
               .. controls (1.08,3.0) and (1.85,2.45) .. (2.45,.85);
  \fill (-1.65,.85) circle (1.5pt);
  \fill (2.45,.85) circle (1.5pt);
  \fill (.30,3.0) circle (1.5pt);
  \draw[densely dotted] (.30,3.0)--(-2.25,3.0);
  \node[left] at (-2.25,3.0) {$t$};
  \node[above] at (.30,3.0) {$(\phi(t),t)$};
  \draw[densely dotted] (-1.65,0)--(-1.65,.85);
  \draw[densely dotted] (2.45,0)--(2.45,.85);
  \node[below] at (-1.65,0) {$y_<$};
  \node[below] at (2.45,0) {$y_>$};
  \draw[densely dotted] (.4,.85)--(.4,-.55);
  \draw[<->] (-1.65,.53)--(.4,.53) node[midway,below] {$r$};
  \draw[<->] (.4,.53)--(2.45,.53) node[midway,below] {$r$};
  \draw[densely dotted] (-.45,.85)--(-.45,-.55);
  \node[below] at (-.45,-.65) {$\phi(h)$};
  \node[below] at (.95,-.65) {$\phi(h)+M$};
  \draw[thin] (.4,-.55)--(.95,-.65);
  \draw[<->] (-.45,-.42)--(.4,-.42) node[midway,above] {$M$};
\end{tikzpicture}
\caption{An arc at base height $h$ and peak height $t$, drawn
schematically with $M>0$. Its endpoint half-width is $r$; its midpoint
is measured relative to $\phi(h)$. Both endpoints move outwards as
the peak rises, giving $|M_r|<1$. The two physical half-orbits form
a periodic orbit precisely when their height-zero endpoint pairs
agree, as proved in Proposition~\ref{prop:cycle-matching}.}
\label{fig:arc-coordinates}
\end{figure}
